% 19-research-foundations-and-principles.tex
% Chapter 10: The Theory of Product and Service Design
% Intellectual Foundations, Eleven Principles, and Conclusion

\section*{Intellectual Foundations and Annotated References}
\addcontentsline{toc}{section}{Intellectual Foundations and Annotated References}

The framework articulated throughout this chapter synthesizes insights across human-centered design, quality engineering, service operations, behavioral ergonomics, and reliability science. These authoritative standards and research traditions provide rigorous conceptual clarity, not a guarantee of commercial success.

\begin{enumerate}[leftmargin=*, label={[\textbf{R\arabic*}]}]
\item \textbf{IDEO, ``Design Thinking at IDEO'':} Explains the interaction of desirability, technical feasibility, and organizational viability, together with iterative human-centered development. Used here as an organizing lens rather than a universal three-variable optimization formula.

\item \textbf{ISO 9241-210:2019, \emph{Human-centred design for interactive systems}:} Requirements and recommendations for understanding users, tasks, environments, design participation, and iterative evaluation. The 2019 edition was reviewed and confirmed in 2025. The standard is specific to interactive systems; its underlying discipline is applied here by analogy to other services.

\item \textbf{Design Council, ``The Double Diamond'' and ``Framework for Innovation'':} The discover-define-develop-deliver representation of divergent and convergent design thinking. It is a process aid, not a rigid law that every startup must perform identical workshops.

\item \textbf{John R. Hauser and Don Clausing (1988), ``The House of Quality'', \emph{Harvard Business Review}:} An influential account of quality-function deployment connecting customer requirements to technical design characteristics across functions. Used as inspiration for lightweight traceability rather than requiring full formal QFD for every startup.

\item \textbf{Noriaki Kano, Nobuhiko Seraku, Fumio Takahashi, and Shin-ichi Tsuji (1984), ``Attractive Quality and Must-Be Quality'':} Useful for understanding why basic expectations and attractive attributes differ; categories can vary by segment and evolve over time.

\item \textbf{Don Norman (2013), \emph{The Design of Everyday Things}, revised and expanded edition:} Important framework for discoverability, affordances, signifiers, feedback, mapping, constraints, conceptual models, and the psychology of everyday interactions.

\item \textbf{Jakob Nielsen (1994; reviewed 2024), ``10 Usability Heuristics for User Interface Design'':} General interaction-design heuristics including visibility of system status, user control, error prevention, recognition over recall, and recovery. These are heuristics, not proof that a specific interface is usable.

\item \textbf{Sarah Gibbons and Alita Kendrick, Nielsen Norman Group (2019), ``Service Blueprinting: Top Questions Answered'':} Clear distinction between customer journey maps and service blueprints, especially customer experience versus frontstage, backstage, and support systems.

\item \textbf{Jakob Nielsen (2012), ``How Many Test Users in a Usability Study?'':} Explains why small iterative qualitative usability studies can be efficient, with major caveats for different studies and quantitative inference. Five test users is a commonly used formative heuristic, not a universal sample size.

\item \textbf{G. Lynn Shostack (1984), ``Designing Services That Deliver'', \emph{Harvard Business Review}:} Foundational treatment of service blueprinting, designing service processes, and anticipating failure points.

\item \textbf{National Institute of Standards and Technology, \emph{AI Risk Management Framework}, version 1.0 (2023):} Voluntary approach to trustworthy AI risks, including validity, reliability, security, accountability, privacy, explainability, and fairness. The framework is under further development; application requires context-specific judgment.

\item \textbf{ISO/IEC 25010:2023, \emph{Systems and software Quality Requirements and Evaluation: Product quality model}:} Specifies a product-quality reference model for ICT/software products consisting of 9 characteristics and subcharacteristics. It does not directly supply a quality model for education or all services.

\item \textbf{Google, \emph{Site Reliability Engineering}, ``Embracing Risk'':} Introduces objective reliability targets and error budgets as a way to balance system reliability and product change. The techniques must be adapted for systems with different failure stakes.

\item \textbf{W. Edwards Deming Institute, ``PDSA Cycle'':} Plan-Do-Study-Act as a method for testing improvement theories, studying results, and adjusting processes. Used here as a philosophy of disciplined iteration, not as a claim that every process change can be experimentally isolated.

\item \textbf{Eric Ries, ``What Is an MVP?'':} Explains that an MVP concerns validated learning with minimal effort, not merely creating a low-feature or poor-quality product.
\end{enumerate}

\subsection*{How the Intellectual Traditions Synthesize}

These diverse traditions form a coherent intellectual continuum. Human-centered design identifies authentic customer needs and environmental contexts; Quality Function Deployment maps those needs to rigorous engineering specifications; Norman and Nielsen eliminate behavioral friction and cognitive overload; Service blueprinting exposes the invisible operational machinery and human handoffs; Reliability engineering and Deming's quality science guarantee consistency under peak stress; and Lean experimentation reduces the capital cost of discovering a working configuration.

No single methodology is sufficient in isolation. A venture that masters interface usability while ignoring operational queues creates delightful experiences it cannot afford to deliver. A venture that optimizes production efficiency while ignoring customer mental models manufactures products nobody wants. The founder's art is the deliberate orchestration of these traditions into a unified, sustainable value-creation system.

\section*{Eleven Principles to Carry Forward}
\addcontentsline{toc}{section}{Eleven Principles to Carry Forward}

\begin{enumerate}[leftmargin=*]
\item \textbf{Design the customer's change of state, not your feature inventory.} A product is valuable only when it reliably makes a consequential job easier, faster, safer, or newly achievable relative to existing alternatives.
\item \textbf{The whole operating system is part of the product.} Onboarding, identity verification, documentation, billing, customer support, handoffs, and error recovery are experienced by the customer as the product itself.
\item \textbf{Complete the job within the promised scope.} A narrow, highly bounded outcome delivered with absolute reliability creates far more trust than an expansive vision riddled with broken dependencies.
\item \textbf{Use requirements that can be independently tested.} Replace vague aspirations like `world-class mentorship' with explicit rubrics, observable value events, turnaround thresholds, and accountable owners.
\item \textbf{Protect the user's real effort.} Eliminate pointless administrative friction and cognitive clutter while fiercely defending the productive intellectual strain and practice essential to developing capability.
\item \textbf{Design frontstage and backstage together.} The visible customer experience is merely the surface reflection of backstage labor, queue discipline, reviewer calibration, and technical infrastructure.
\item \textbf{Reliability is an end-to-end series system.} A sequence of individually competent components can still produce an untrustworthy customer outcome if handoffs are unmonitored and error propagation is uncontained.
\item \textbf{Minimum viable does not mean minimum responsible.} Bounding early experimental scope must never serve as an excuse for reckless safety, deceptive promises, or omitting the core value mechanism.
\item \textbf{Iteration must alter decisions, not merely create activity.} Measure observable value events, failure severities, and fully loaded delivery costs rather than celebrating vanity engagement metrics.
\item \textbf{Customer value and operating viability must compound together.} A beloved customer service that relies on permanent, uncompensated founder overtime is an unsustainable illusion, not a business.
\item \textbf{Build institutional capabilities that outlive the founder.} Every completed engagement should leave behind reusable instructional assets, calibrated rubrics, codified operating procedures, and empirical evidence.
\end{enumerate}

\subsection*{Final Thought: A Promise That Becomes a Capability}

The early founder begins with an imaginative vision of what a customer might desire. They assemble a prototype, pitch an inspiring demonstration, and hope that the market recognizes their ingenuity. In the beginning, this creative leap is indispensable. Without imagination, nothing new enters the world.

However, an enduring enterprise does not emerge from imagination alone. It emerges when the founder can reliably coordinate people, software, physical materials, operational procedures, human judgment, and economic incentives so that the promised transformation occurs under ordinary, messy operating conditions. They understand how to detect defects, how to execute graceful recovery, what the customer must contribute, what the offering can truthfully claim, and what each delivery costs in human hours and capital.

A serious academic institution makes this distinction unmistakable. It is not enough to possess video recordings and lecture slides. The institution must engineer the precise operating conditions under which eligible learners can undertake demanding intellectual work, receive unsparing expert critique, correct profound misconceptions, demonstrate independent mastery, and retain durable capabilities long after the course concludes. The institution's ultimate achievement is not instructional production; it is \textbf{reproducible human capability}.

The first time a founder produces this transformation, it demands heroic personal exertion. The second time, codified operating procedures and calibrated rubrics alleviate the burden. By the tenth cohort, an operating team executes the mechanism with calibrated consistency. Eventually, the founder's personal brilliance becomes an institutional capability: a promise made reliably true.

\textbf{That is what a product really is: a promise made operational.}
